\documentclass[12pt]{article}
\usepackage{amsmath}
\usepackage{amsfonts}
\usepackage{amssymb}
\usepackage{geometry}
\usepackage{tikz}
\geometry{a4paper, margin=1in}
\pagenumbering{gobble}


\setlength\parindent{0pt}

\title{02YMECH - Homework 10}
\date{Assigned for the week of Nov 24, 2025}
\author{}
\begin{document}
\maketitle

\section*{Question}

\begin{enumerate}
    \item Check the validity of the phrase ``apple doesn't fall far from the tree'' for an apple tree at the equator after considering the Coriolis force. Consider that the apple falls from the height $h = 4$~m. Neglect the air resistance.
    
    \underline{Hint}: Neglect the effect of the horizontal velocity caused by the Coriolis force: If the falling apple starts with zero horizontal velocity, then its horizontal velocity appears only because of the Coriolis force itself. Strictly speaking, once this horizontal component exists, it should generate a second-order Coriolis correction (because now velocity has a horizontal piece). That correction is extremely small, neglect it.
    
    \item A particle of mass $m$ moves without friction inside a smooth, rigid, straight pipe that is fixed on a horizontal rotating platform. The pipe is oriented radially outward from the axis of rotation. The platform rotates with a constant angular velocity $\Omega$. Initially, the particle is held at radius $r_0$ with zero velocity relative to the pipe and then released.

\begin{enumerate}
    \item In the rotating frame of the platform, write down the equation of motion for the particle, including all inertial forces.
    \item Solve for $r(t)$.
    \item Determine the time it takes for the particle to reach the end of the pipe at radius $R > r_0$.
    \item Explain physically why the particle accelerates, even though its velocity is initially zero in the rotating frame.
\end{enumerate}
    


\end{enumerate}

\end{document}
